%قبل از شروع مسئله ترجیحاً چیزی ننویسید. اما اگر خواستید بنویسید، در قالب لیستی (فقرات و شمارش) بنویسید.
%برای اضافه کردن نمره، بعدی از شروع‌مسئله، [x] را اضافه کنید و به جای x  نمره مورد نظر را قرار دهید. برای زیرمسئله‌ها هم به همین صورت است.
%\شروع‌مسئله[10]
\شروع‌مسئله [20]
\متن‌سیاه{نابه‌جایی با AVL}

در آرایۀ $A = (a_1, a_2, ..., a_n)$ به زوج $i$ و $j$ که $i > j$ و $a_i < a_j$ یک نابه‌جایی می‌گویند. یک آرایۀ $n$ عضوی به شما داده می‌شود و شما باید برای هر زیردنبالۀ $k$تایی متوالی آن تعداد نابه‌جایی‌ها را به دست آورید و سپس مجموع این اعداد را حساب کنید. برای حساب کردن این مجموع شما باید از درخت AVL استفاده کنید.

\متن‌سیاه{ورودی:}
در خط اول به ترتیب اعداد $n$ و $k$ به شما داده می‌شود. در خط بعد $n$ عدد به ترتیب به شما داده می‌شود.

\متن‌سیاه{خروجی:}
خروجی شما باید مجموع تعداد نابه‌جایی‌ها به پیمانه‌ی $10^9 + 7$ باشد.

\متن‌سیاه{نکات:}
در صورتی که از درخت AVL در پیاده‌سازی خود استفاده نکنید نمره‌ای به شما تعلق نخواهد گرفت.

\متن‌سیاه{ورودی و خروجی‌های نمونه:}
\begin{latin}
	\begin{tabularx}{\textwidth}{|X | X|}
		\hline
		\textbf{input} &
		\textbf{output}\\
		\hline
		
		6 2 &
		\\
		
		1 2 3 4 5 6 &
		0\\
		\hline
		4 3 &
		\\
		
		1 6 7 2 &
		2\\
		\hline
		10 5&
		 \\
		
		15 51 44 44 76 50 29 88 48 50 &
		25\\
		\hline
		8 4 &
		 \\
		
		12 3 14 8 15 1 4 22&
		14\\
		
		\hline
		
	\end{tabularx}
\end{latin}



\solution{
درخت \متن‌لاتین{AVL}ای که برای حل این سؤال استفاده می‌کنیم، یک درخت مرتبه آماری با قابلیت درج و حذف است (مشابه آن‌چه برای محاسبۀ مرتبه (rank) در جلسۀ ششم درس مطرح شد). این درخت برای هر رأس، تعداد رئوس زیردرخت آن رأس را ذخیره می‌کند و با استفاده از آن جایگاه هر عنصر را در آرایۀ مرتب شده به دست می‌آورد. ابتدا با $k$ عنصر اول این درخت را می‌سازیم و تعداد نابه‌جایی‌های آن را حساب می‌کنیم. تعداد نابه‌جایی‌هایی که هر عنصر با عناصر قبل از آن ایجاد می‌کند برابر تعداد عناصری است که هنگام درج آن عنصر در درخت از آن بزرگ‌تر هستند. حال می‌خواهیم تعداد نابه‌جایی‌های $k$ عنصر بعدی را پیدا کنیم. ابتدا عنصر اول $k$تایی قبلی را حذف می‌کنیم و تعداد نابه‌جایی‌هایی که ایجاد می‌کرد را از تعداد نابه‌جایی‌های $k$تایی کم می‌کنیم تا تعداد نابه‌جایی‌های $k-1$ تایی به دست آید. سپس عنصر جدید را درج می‌کنیم و تعداد نابه‌جایی‌هایی که ایجاد می‌کند را به تعداد نابه‌جایی‌های $k-1$تایی اضافه می‌کنیم تا تعداد نابه‌جایی‌های $k$ تایی جدید به دست آید. به این صورت می‌توان جمع تعداد نابه‌جایی‌های $k$تایی‌ها را در زمان $O(n\log k)$ به دست آورد. کد ضمیمه شده برای این سوال پاسخ ارسالی آقای وحید زهتاب می‌باشد که در فایل \متن‌تایپ{avl.py} موجود است.


}




%%%%%%%%%%%%%%%%%%%%%%%%%%%%%%%%%%%%%%%%%%%%%%%%%%%%%
% برخی دستورات مفید برای تمرین‌ها. %
%%%%%%%%%%%%%%%%%%%%%%%%%%%%%%%%%%%%%%%%%%%%%%%%%%%%%

% فرمول‌های نمونه:
%$f(x) = \Theta(x)$
%$T(x, y) \leq \log(x) + 2^y + \binom{2n}{n}$
% {\tt code\_function}


% لیست موارد بدون شماره‌گذاری:
%\شروع{فقرات}
%    \فقره مورد اول 
%        \شروع{فقرات}
%            \فقره مورد اول داخلی 
%                \شروع{فقرات}
%                    \فقره اولین مورد خیلی داخلی 
%                    \فقره دومین مورد خیلی داخلی 
%                \پایان{فقرات}
%            \فقره مورد دوم داخلی 
%        \پایان{فقرات}
%    \فقره مورد دوم 
%\پایان{فقرات}

% لیست موارد شماره‌گذاری شده:
%\شروع{شمارش}
%    \فقره یک مورد 
%    \فقره یک مورد دیگر 
%    \فقره باز هم یک مورد دیگر
%\پایان{شمارش}

% برای نوشتن یک فرمول چند خطی که یک شماره بخوره. برای حالت‌های دیگر به سایت زیر مراجعه کنید:
% https://www.overleaf.com/learn/latex/Aligning_equations_with_amsmath
%\begin{align} 
%    \begin{split}
%        (x+y)^3 &= (x+y)^2(x+y) \\
%                &= (x^2+2xy+y^2)(x+y) \\
%                &= (x^3+2x^2y+xy^2) + (x^2y+2xy^2+y^3) \\
%                &= x^3+3x^2y+3xy^2+y^3
%    \end{split}                                 
%\end{align}

% نوشتن ماتریس و آرایه دوبعدی
%\begin{align}
%    A = 
%    \begin{bmatrix}
%        A_{11} & A_{21} \\
%        A_{21} & A_{22}
%    \end{bmatrix}
%\end{align}

% اضافه کردن تصویر با شرح. برای اضافه کردن تصویر بدون شرح، فقط دستور درج‌تصویر را استفاده کنید.
%\شروع{شکل}[h!]
%\تنظیم‌ازوسط
%  \درج‌تصویر[پهنا=0.5\پهنای‌متن]{img.jpg}
%\شرح{شرح تصویر}
%\برچسب{تصویر:۱}
%\پایان{شکل}
%\پایان{زیرمسائل}

%کد زیر را ببینید.
%%\lstinputlisting{source_filename.py}
%\begin{latin}
%\begin{lstlisting}
%# Import the modules
%import sys
%import random
%
%ans = True
%
%while ans:
%question = raw_input("Ask the magic 8 ball a question: (press enter to quit) ")
%
%answers = random.randint(1,8)
%
%if question == "":
%sys.exit()
%
%elif answers == 1:
%print "It is certain"
%
%elif answers == 2:
%print "Outlook good"
%
%elif answers == 3:
%print "You may rely on it"
%
%elif answers == 4:
%print "Ask again later"
%
%elif answers == 5:
%print "Concentrate and ask again"
%
%elif answers == 6:
%print "Reply hazy, try again"
%
%elif answers == 7:
%print "My reply is no"
%
%elif answers == 8:
%print "My sources say no"
%\end{lstlisting}
%\end{latin}
